\section{襄国、成都与建康之间：三三〇至三四九年的多重重组}
\label{sec:eastern-jin-post-zhao-chenghan}

三二九年前赵灭亡后，石勒在次年于襄国称帝，后赵成为北方最强的政权之一。称帝与改元可靠，官制仿效到什么程度、羯胡与汉人官僚如何共处、地方命令是否真正执行到各城镇，不能由《晋书》〈石勒载记〉的帝国叙事直接推出。后赵能调动河北的资源，也意味着黄河流域的城郭、军队与家庭必须承担新的征发与秩序。

\subsection{继承危机没有只发生在建康}

三三三年石勒死，太子石弘继位，石虎掌握军政；次年石虎废杀石弘，自立为天王。遗诏、石弘的自主空间和石虎取得支持的过程都有争议，废立与权力转移却无可怀疑。后赵的高压扩张暂时得到更集中的指挥，也把继承危机压入宗室、军队和宫廷之间。三四九年石虎死后，诸子与养孙围绕邺城、军队和皇位相争，说明先前的集中并没有建立能够自我约束的继承程序。

西南的成汉走过相似而不相同的道路。李雄三三四年去世，李班继位后旋即被李期夺权；三三八年李寿再废李期，改国号为汉。成都宫廷中谁支持哪一位宗室、各次废立的罪状是否可靠，不能全凭后出的载记裁断；频繁改号与政变确实削弱了政权的正当性和应变能力。成汉并非因“偏远”而天然脆弱，它需要在巴蜀的官僚、道路和军队之间维持协调，而继承冲突正在消耗这种协调。

\subsection{辽西、丸都与区域的道路}

三三五年慕容皝在辽西称燕王，次年击败段辽并整合部众；三三八年石虎进攻前燕未能摧毁其核心地区。战场细节、部众数量和所谓降附者的迁徙路线有争议，但辽西通道、棘城周边地形和可动员的人口，使前燕得以在后赵的压力下继续发展。三四二年慕容皝进攻高句丽都城丸都，远征和城邑受创可由《晋书》、后出的《三国史记》及《资治通鉴》交叉辨认；俘获人数、王室损失与双方叙事的差异不应被抹平。

这些北方和东北的战事并不是建康能够直接指挥的边缘故事。它们不断改变流民、使节、军镇和地方集团可以选择的道路，也决定东晋北伐面对的是哪一种政权组合。东晋的北方目标因而不能只写成恢复洛阳的愿望，还要看见每一条从江淮通往关中、河北和辽东的路径，都由不同的城郭和人际网络控制。

\subsection{墓志与巴蜀的再接合}

三四一年王兴之夫妇墓志下葬，南京象山的发掘为建康周边侨姓官员家庭、婚姻和官职提供了带纪年的物证。墓志可以校正姓名、日期和某些仕宦关系，却是家族的公开自我呈现；追赠官、世系和道德语言不能替代更广泛的社会统计。它的价值正在于提示我们，东晋的政治并不只存在于帝纪中的王、桓等人，也存在于迁徙家庭如何安置、婚配和书写身份的地方网络。

三四六年桓温自荆州西征成汉，次年攻入成都，李势降，东晋恢复益州控制。出师、主要路线与成都降服可靠，兵数、补给、战后安置及蜀地整合的速度均有争议。东晋取得上游的税源和战略纵深，也显著抬高了桓温的个人声望；巴蜀重新进入建康的任命网络，却不等于当地社会立即按江南中枢的意图运转。成汉覆灭与后赵继承危机相互并行，使三四九年以后的北方又一次进入更剧烈的重组。

\subsection{材料的边界}

《晋书》〈石勒载记〉〈石虎载记〉、慕容与李氏载记，《华阳国志》以及《资治通鉴》共同保存这一时期的政治主线；十六国史多已散佚，后代辑佚文字能补国号和地名，却不能被当作完整原史。王兴之夫妇墓志和发掘报告提供独立的地域与纪年锚点，但不能代表建康的全部人口。由此可见的，是襄国、成都、棘城和建康之间并无一条单向的强弱线，而是多种继承、迁徙和军政网络同时改变的历史空间。
